\documentclass[../Main.tex]{subfiles}

\begin{document}

\section{Reflexion} \label{sec:Reflexion}

\subsection{Reflexion des Arbeitsprozesses}

Unser Projekt dauerte etwa ein Jahr und durchlief mehrere Phasen, die wir im Folgenden der Reihe nach reflektieren.

Am Anfang hatten wir noch ganz andere Vorstellungen vom Endprodukt und stocherten in alle Bereiche hinein: Wir bauten Spielsysteme, von denen wir ohne Recherche annahmen, dass wir sie brauchen, überlegten, wie wir die Erkenntnisfrage beantworten wollen, und lernten zu recherchieren. Bis zur Zwischenpräsentation am 17.~Juni~2026 blieben uns noch vier Monate. Das wirkt planlos, war rückblickend aber ein sinnvolles Auskundschaften, zumal noch wenig auf dem Spiel stand: Nichts musste verworfen werden, und die Arbeitsschritte hingen kaum voneinander ab. Dass es gut ging, war zwar eher Zufall als Planung. Trotzdem würden wir es wieder so angehen, nur im Bewusstsein, in dieser Phase noch eine Narrenfrist zu haben.

Bis zur Zwischenpräsentation stand die Erkenntnisfrage im Hintergrund, denn die Playtests setzten eine technische Grundlage voraus, die erst entstehen musste: den Spielercharakter, die Gegner und alles, was das Spiel überhaupt zum Laufen bringt. Die Recherche kam dabei zu kurz, rückblickend war das aber die richtige Entscheidung. Dieser Teil machte uns am meisten Freude, und erst durch die eigenen Systeme wussten wir, worauf wir in anderen Spielen achten mussten. Künftig würden wir zuerst recherchieren, wenn unsere Vorstellungen konkret sind und das nötige Vorwissen vorhanden ist, und sonst zuerst eigene Erfahrungen sammeln.

Vor der Zwischenpräsentation mussten wir Features, die für die Erkenntnisfrage wenig relevant waren, streichen oder eingrenzen und zusätzlich einen begründeten Playtest durchführen. Wir erreichten alle vertraglich festgehaltenen Ziele, aber nur sehr knapp, vor allem weil wir bei der Qualität kaum Kompromisse eingehen wollten. Da die Zwischenbewertung nicht in die Note der Maturitätsarbeit einfliesst, haben wir die Chance verpasst zu erfahren, wie 80 bis 85 Prozent unseres Anspruchs bewertet worden wären. Dafür sammelten wir Erfahrung mit Crunch, die uns am Ende half, mit genügend Zeit und kaum Crunch 95 bis 100 Prozent herauszuholen.

Die Kritik der Zwischenbewertung betraf vor allem die begleitenden Texte. Da wir die bemängelten Punkte zunächst nicht verinnerlicht hatten, fragten wir bei unserer Betreuungsperson so lange nach, bis wir sie verstanden.

Danach erweiterten wir den Projektumfang, da die Abgabe noch weit entfernt war. Gleichzeitig sank unser Tempo: Bei der Grösse des Projekts machten kleine Neuerungen kaum noch einen Unterschied, und die Pflege des Projekts sowie angehäufte, weniger spannende Aufgaben setzten uns zu, auch wenn unser Eifer nicht erloschen war.

\begin{itemize}[leftmargin=1.5cm]
    \item \textcolor{red}{\textbf{1. Ausgangslage und Plan:} Erste Commits im Juni 2025; der Vertrag sah vor: Zwischenprodukt 10.6., Zwischenpräsentation 17.6., Analyse der Referenzspiele Anfang Juli, Spielsysteme Mitte August, Doku und Fragebögen Ende September, Playtests Mitte Oktober, Verbesserungen Ende Oktober.}
    \item \textcolor{red}{\textbf{2. Verlauf in Phasen} (aus Git und Journal): Hauptentwicklung des Prototyps Februar bis Juni 2026; Juni: README und Architekturdiagramm fürs Zwischenprodukt; 16.--20.7.: Kapitel 1--4 und die \game{NieR}-Analyse als Entwurf; Juli bis September kaum Arbeit am Spiel (Sommerpause?); Oktober: Level überarbeitet, Playtests vorbereitet und durchgeführt; November: Begleitbericht fertiggestellt.}
    \item \textcolor{red}{\textbf{3. Scope-Entscheide:} 20.5.: Vertrag bereinigt, Projektumfang reduziert, Erkenntnisfrage festgelegt -- was wurde gestrichen? Später: Bonus-Inhalte weggelassen, technische Dokumentation ausgelagert, README gestrichen, Recherche in Kapitel~\ref{sec:Leveldesign} auf Rat der Betreuungsperson auf das für die Erkenntnisfrage Nötige gekürzt.}
    \item \textcolor{red}{\textbf{4. Plan und tatsächlicher Verlauf:} Playtests 26.10.--2.11.\ statt Mitte Oktober, Verbesserungen erst im Bericht statt Ende Oktober. Warum?}
    \item \textcolor{red}{\textbf{5. Wichtige Wendepunkte} (Journal 24./25.10., 13~h): Obwohl der Prototyp technisch fertig war, machten die Kämpfe bei eigenen Tests wenig Spass. Ihr habt die Encounters gekürzt, die Gegner aggressiver gemacht und Rechercheergebnisse nachgeholt, die noch nicht umgesetzt waren. Eure eigene Erkenntnis dort: \enquote{Die technische Fertigstellung allein garantiert noch kein gutes Spielerlebnis} -- und weil ihr das Level so oft gespielt habt, braucht es externe Playtests. Das ist ein starkes Reflexionsbeispiel.}
    \item \textcolor{red}{\textbf{6. Rückmeldungen der Zwischenbewertung und was ihr daraus gemacht habt:} Zeitaufwand konsequenter dokumentieren, Reflexionen vertiefen, Entscheidungen mit Alternativen begründen, KI vollständig dokumentieren, Planung mit Arbeitspaketen, Bericht und Code verknüpfen.}
    \item \textcolor{red}{\textbf{7. Was gut funktioniert hat:} z.\,B. iteratives Vorgehen, Git und GitHub, Playtest-Vorgehen -- jeweils mit Grund.}
    \item \textcolor{red}{\textbf{8. Methodik im Rückblick:} kurz, Einzelheiten in den Limitationen (Abschnitt~\ref{sec:Limitationen}).}
    \item \textcolor{red}{\textbf{9. Was ihr anders machen würdet}, z.\,B. früher extern testen, Rechercheergebnisse von Anfang an mit dem Level abgleichen, Bericht laufend statt am Schluss schreiben.}
    \item \textcolor{red}{\textbf{10. Verweis aufs Arbeitsjournal}, wo der Verlauf im Detail dokumentiert ist.}
\end{itemize}

\subsection{Gruppenarbeit und Aufgabenteilung}

\begin{itemize}[leftmargin=1.5cm]
    \item \textcolor{red}{\textbf{1. Wie ihr die Arbeit aufgeteilt habt.} Laut Zwischenbewertung: Rodin stärker bei Animation, UI, Bericht und README; Zhile stärker bei Architektur, Bewegung, State Machines und Kampf-Grundlagen. Das deckt sich mit Git: Zhile vor allem \code{Entities/Player}, \code{Common} (Utility, Tools), \code{Core}, Kamera; Rodin vor allem \code{Interfaces} und Menüs sowie die TrenchBroom-Karten des Levels.}
    \item \textcolor{red}{\textbf{2. Tabelle Aufgabenteilung und Eigenleistung} für Prototyp und Bericht (Rodin / Zhile / gemeinsam). Zum Bericht: Kapitel 1--8 und die \game{NieR}-Analyse Rodin, die \game{Ninja-Gaiden}-Analyse Zhile, README und Architekturdiagramm (Juni) vor allem Zhile. Bitte prüfen und ergänzen (Playtests, Leveldesign, Audio, \ldots).}
    \item \textcolor{red}{\textbf{3. Zusammenarbeit und Kommunikation:} Wie habt ihr euch abgesprochen (Treffen, Discord, GitHub)? Wie habt ihr Entscheidungen getroffen?}
    \item \textcolor{red}{\textbf{4. Veränderung im Verlauf:} z.\,B. ab Oktober Rodin vor allem Level, Playtests und Bericht, Zhile weiter am Spiel (laut Git im Oktober sehr aktiv). Wie kam es dazu?}
    \item \textcolor{red}{\textbf{5. Fazit:} Was lief gut, was war schwierig, was würdet ihr in einer Gruppenarbeit anders machen? Hinweis: Commit-Zahlen sagen wenig über den Aufwand aus, deshalb nicht als Massstab verwenden.}
\end{itemize}

\subsection{Verwendung von KI} \label{sec:ReflexionKI}

\begin{itemize}[leftmargin=1.5cm]
    \item \textcolor{red}{\textbf{1. Welche Tools wofür:} Code (Treppen-Werkzeug und Animations-Hilfsskript im Frühling, Audiosystem im Oktober), Begleitbericht (Korrektur, Struktur, Rückmeldungen). Details in der KI-Deklaration (Abschnitt~\ref{sec:KIDeklaration}).}
    \item \textcolor{red}{\textbf{2. Was geholfen hat:} z.\,B. Widersprüche zwischen Kapiteln finden, unbelegte Aussagen markieren, \LaTeX-Formatierung, Fehlersuche im Code.}
    \item \textcolor{red}{\textbf{3. Grenzen und Probleme:} Die KI ging mehrmals von falschen Annahmen aus, die ihr korrigieren musstet, z.\,B. dass Themenfindung mit KI gemacht wurde, dass ein Tutorial geplant war, welche Räume Brückenräume sind, oder dass in den Protokollen ein anderer Archer gemeint war.}
    \item \textcolor{red}{\textbf{4. Wie ihr die Ergebnisse überprüft habt:} Abgleich mit Quellen, Playtest-Protokollen und dem Spiel selbst; Vorschläge nicht ungeprüft übernommen, einige bewusst verworfen.}
    \item \textcolor{red}{\textbf{5. Eigenanteil sichern:} Interpretation und Urteile selbst festgelegt; von der KI formulierte Stellen markiert und in eigenen Worten neu geschrieben; Regel \enquote{nicht wortwörtlich übernehmen} aus dem Merkblatt.}
    \item \textcolor{red}{\textbf{6. Fazit:} Würdet ihr KI wieder so einsetzen? Was würdet ihr anders machen, z.\,B. KI-Nutzung von Anfang an laufend im Journal dokumentieren?}
\end{itemize}

\subsection{Persönliche Erkenntnisse}

\subsubsection{Rodin Arslan}

\begin{itemize}[leftmargin=1.5cm]
    \item \textcolor{red}{\textbf{1. Fachlich:} z.\,B. Animation (u.\,a. Hängeanimation), UI und Menüs, Leveldesign mit TrenchBroom, Gegner- und Encounter-Design, Playtesting, \LaTeX\ und BibTeX.}
    \item \textcolor{red}{\textbf{2. Methodisch:} Spielanalyse mit dem MDA-Framework, Recherche und Zitieren, Playtests planen und auswerten, Projektplanung.}
    \item \textcolor{red}{\textbf{3. Persönlich:} z.\,B. dass ein technisch fertiges Spiel noch nicht gut spielbar ist (Journal 24./25.10.), Zeitmanagement in der Schlussphase, Umgang mit Rückmeldungen.}
    \item \textcolor{red}{\textbf{4. Was du mitnimmst} für die Zukunft.}
\end{itemize}

\subsubsection{Zhile Zhu}

\textcolor{red}{ZHILE DAS IST FÜR DICH}

\end{document}
